% Options for packages loaded elsewhere
\PassOptionsToPackage{unicode}{hyperref}
\PassOptionsToPackage{hyphens}{url}
\documentclass[
  ignorenonframetext,
]{beamer}
\newif\ifbibliography
\usepackage{pgfpages}
\setbeamertemplate{caption}[numbered]
\setbeamertemplate{caption label separator}{: }
\setbeamercolor{caption name}{fg=normal text.fg}
\beamertemplatenavigationsymbolsempty
% remove section numbering
\setbeamertemplate{part page}{
  \centering
  \begin{beamercolorbox}[sep=16pt,center]{part title}
    \usebeamerfont{part title}\insertpart\par
  \end{beamercolorbox}
}
\setbeamertemplate{section page}{
  \centering
  \begin{beamercolorbox}[sep=12pt,center]{section title}
    \usebeamerfont{section title}\insertsection\par
  \end{beamercolorbox}
}
\setbeamertemplate{subsection page}{
  \centering
  \begin{beamercolorbox}[sep=8pt,center]{subsection title}
    \usebeamerfont{subsection title}\insertsubsection\par
  \end{beamercolorbox}
}
% Prevent slide breaks in the middle of a paragraph
\widowpenalties 1 10000
\raggedbottom
\AtBeginPart{
  \frame{\partpage}
}
\AtBeginSection{
  \ifbibliography
  \else
    \frame{\sectionpage}
  \fi
}
\AtBeginSubsection{
  \frame{\subsectionpage}
}
\usepackage{iftex}
\ifPDFTeX
  \usepackage[T1]{fontenc}
  \usepackage[utf8]{inputenc}
  \usepackage{textcomp} % provide euro and other symbols
\else % if luatex or xetex
  \usepackage{unicode-math} % this also loads fontspec
  \defaultfontfeatures{Scale=MatchLowercase}
  \defaultfontfeatures[\rmfamily]{Ligatures=TeX,Scale=1}
\fi
\usepackage{lmodern}
\ifPDFTeX\else
  % xetex/luatex font selection
\fi
% Use upquote if available, for straight quotes in verbatim environments
\IfFileExists{upquote.sty}{\usepackage{upquote}}{}
\IfFileExists{microtype.sty}{% use microtype if available
  \usepackage[]{microtype}
  \UseMicrotypeSet[protrusion]{basicmath} % disable protrusion for tt fonts
}{}
\makeatletter
\@ifundefined{KOMAClassName}{% if non-KOMA class
  \IfFileExists{parskip.sty}{%
    \usepackage{parskip}
  }{% else
    \setlength{\parindent}{0pt}
    \setlength{\parskip}{6pt plus 2pt minus 1pt}}
}{% if KOMA class
  \KOMAoptions{parskip=half}}
\makeatother
\usepackage{longtable,booktabs,array}
\newcounter{none} % for unnumbered tables
\usepackage{calc} % for calculating minipage widths
\usepackage{caption}
% Make caption package work with longtable
\makeatletter
\def\fnum@table{\tablename~\thetable}
\makeatother
\setlength{\emergencystretch}{3em} % prevent overfull lines
\providecommand{\tightlist}{%
  \setlength{\itemsep}{0pt}\setlength{\parskip}{0pt}}
% Auto-generated by infrastructure/rendering/_pdf_latex_helpers.py::extract_math_font_preamble.
% Minimal header passed to Pandoc via -H so Beamer slides pick up the same math font
% as the combined-PDF path. Adjust the manuscript's preamble.md to change the font globally.
\usepackage{unicode-math}
\setmathfont{latinmodern-math.otf}

% Slide layout defaults for warning-clean scientific decks.
\usepackage{etoolbox}
\IfFileExists{xurl.sty}{\usepackage{xurl}}{}
\IfFileExists{seqsplit.sty}{\usepackage{seqsplit}}{\newcommand{\seqsplit}[1]{#1}}
\protected\def\breakseq#1{\seqsplit{#1}}
\protected\def\breaktt#1{\begingroup\ttfamily\seqsplit{#1}\endgroup}
\setlength{\emergencystretch}{6em}
\tolerance=5000
\hbadness=10000
\hfuzz=1pt
\setlength{\tabcolsep}{2pt}
\AtBeginEnvironment{longtable}{\tiny\renewcommand{\arraystretch}{0.86}\setlength{\tabcolsep}{1pt}}
\AtBeginEnvironment{tabular}{\tiny\renewcommand{\arraystretch}{0.86}\setlength{\tabcolsep}{1pt}}

% Natbib and cross-reference fallbacks — slides don't load natbib
% or cleveref, but combined-PDF manuscript prose may emit these
% commands. The fallback renders citations as a bracketed key list
% and cross-references as detokenized labels so slides don't fail on
% undefined control sequences or raw underscores. \providecommand is
% a no-op if packages load later.
\providecommand{\citep}[1]{[#1]}
\providecommand{\citet}[1]{#1}
\providecommand{\citealp}[1]{#1}
\providecommand{\citeauthor}[1]{#1}
\providecommand{\citeyear}[1]{#1}
\providecommand{\cref}[1]{\texttt{\detokenize{#1}}}
\providecommand{\Cref}[1]{\texttt{\detokenize{#1}}}
\usepackage{bookmark}
\IfFileExists{xurl.sty}{\usepackage{xurl}}{} % add URL line breaks if available
\urlstyle{same}
% fallback for those not using the hyperref driver hyperxmp:
\makeatletter
\@ifundefined{xmpquote}{\newcommand{\xmpquote}[1]{#1}}{}
\makeatother
\hypersetup{
  hidelinks,
  pdfcreator={LaTeX via pandoc}}

\author{\texorpdfstring{}{}}
\date{}

\begin{document}

\section{Experimental Setup}\label{sec:experimental_setup}

\begin{frame}[allowframebreaks]{Experimental Setup}
This section details the dataset, schema, and software environment used
to produce the results. The exemplar deliberately avoids manuscript
token injection: concrete numbers are reproduced by running the analysis
script, and the figures are regenerated from tested code, so nothing
here can silently drift from a hardcoded value.
\end{frame}

\begin{frame}[fragile,allowframebreaks]{Dataset}
\protect\phantomsection\label{dataset}
The analysis uses a shipped, deterministic CSV fixture
(\breaktt{data/measurements.csv}) generated once with a fixed NumPy seed.
It contains 120 subject records with the following columns:

{\def\LTcaptype{none} % do not increment counter
\begin{longtable}[]{@{}lll@{}}
\toprule\noalign{}
Column & Role & Type \\
\midrule\noalign{}
\endhead
\texttt{subject\_id} & per-row identifier & string \\
\texttt{group} & categorical group (\texttt{alpha}, \texttt{beta},
\texttt{gamma}) & string \\
\texttt{height\_cm} & numeric feature & float \\
\texttt{weight\_kg} & numeric feature & float \\
\texttt{resting\_hr\_bpm} & numeric feature & float \\
\bottomrule\noalign{}
\end{longtable}
}

These roles are declared once in
\breaktt{src/eda/dataset.py::DatasetSchema}, which the statistics,
correlation, and figure functions consult. The generating process makes
weight depend positively on height (a strong correlation) while resting
heart rate is only weakly related, and a few numeric cells are left
blank to exercise the missing-data path.
\end{frame}

\begin{frame}[fragile,allowframebreaks]{Analysis conditions}
\protect\phantomsection\label{analysis-conditions}
The experiment overlay (\breaktt{experiment\_plan.yaml}) declares three
conditions:

\begin{itemize}
\tightlist
\item
  \textbf{raw\_dataset} (reference) --- as loaded, with missing cells as
  \texttt{NaN}.
\item
  \textbf{cleaned\_dataset} (proposed) --- rows with any missing numeric
  feature dropped via \texttt{clean\_dataset()}, which reports the count
  removed.
\item
  \textbf{normalized\_dataset} (variant) --- the cleaned dataset with
  numeric columns z-score standardized by \breaktt{normalize\_numeric()}
  for cross-feature comparison.
\end{itemize}

The primary descriptive lens is the pairwise Pearson correlation among
the numeric features.
\end{frame}

\begin{frame}[fragile,allowframebreaks]{Computational environment}
\protect\phantomsection\label{computational-environment}
\begin{itemize}
\tightlist
\item
  \textbf{Language}: Python (see root \texttt{pyproject.toml} for the
  supported range).
\item
  \textbf{Core dependencies}: \texttt{numpy}, \texttt{pandas},
  \texttt{matplotlib} (declared in
  \breaktt{domain\_profile.yaml::required\_packages}).
\item
  \textbf{Headless plotting}: the analysis script sets
  \texttt{MPLBACKEND=Agg} before importing matplotlib.
\end{itemize}
\end{frame}

\begin{frame}[fragile,allowframebreaks]{Pipeline ordering}
\protect\phantomsection\label{pipeline-ordering}
The typical analysis order is:

\begin{enumerate}
\tightlist
\item
  \breaktt{scripts/eda\_analysis.py} --- loads and cleans the dataset,
  then writes \breaktt{output/figures/*.png} and
  \breaktt{output/data/summary\_statistics.csv}, printing each output
  path for manifest collection.
\item
  PDF rendering reads \texttt{manuscript/*.md} and \texttt{config.yaml}
  so figure paths and prose match the analysis that just completed.
\end{enumerate}
\end{frame}

\begin{frame}[fragile,allowframebreaks]{Relation to figures}
\protect\phantomsection\label{relation-to-figures}
{\def\LTcaptype{none} % do not increment counter
\begin{longtable}[]{@{}
  >{\raggedright\arraybackslash}p{(\linewidth - 4\tabcolsep) * \real{0.3333}}
  >{\raggedright\arraybackslash}p{(\linewidth - 4\tabcolsep) * \real{0.3333}}
  >{\raggedright\arraybackslash}p{(\linewidth - 4\tabcolsep) * \real{0.3333}}@{}}
\toprule\noalign{}
\begin{minipage}[b]{\linewidth}\raggedright
Figure (\breakseq{[@sec:results]})
\end{minipage} & \begin{minipage}[b]{\linewidth}\raggedright
Figure-data preparer (\breaktt{src/eda/figures.py})
\end{minipage} & \begin{minipage}[b]{\linewidth}\raggedright
Primary inputs
\end{minipage} \\
\midrule\noalign{}
\endhead
Height histogram & \breaktt{histogram\_data()} & \texttt{height\_cm}
column, 10 bins \\
Rows per group & \breaktt{group\_count\_data()} & \texttt{group}
column \\
Correlation heatmap & \breaktt{correlation\_heatmap\_data()} & all
numeric features \\
\bottomrule\noalign{}
\end{longtable}
}

This table is descriptive documentation only; it is not executed as code
during the build.
\end{frame}

\end{document}
